\section{Further research questions and concrete next targets}
\label{sec:further}

The most useful continuations seek exact identities and reductions.
The questions below distinguish a representation problem, a reduction
problem in a prescribed constant space, and an arithmetic independence
problem; solving the first does not automatically solve the other two.

\subsection{The frozen Gaussian relations at weights seven and nine}

The immediate named-conjecture target remains the exact vanishing of
the current S6 and S8 residuals in
\eqref{eq:source-current-S6}--\eqref{eq:source-current-S8}.
The mixed endpoint $\Li_{p,1}(i,-1)$ is essential.  A useful next result
would be a functional identity involving both endpoints that is valid
before specialization to a fixed weight.  Its specialization could
then be reduced by exact shuffle and stuffle relations to the frozen
baskets.  For S6, a candidate argument can be tested against the
source's separating functional: an argument confined to the already
insufficient restricted presentation cannot suffice.

An exact residual derivation should record the initial analytic domain,
the admissible endpoint continuation, every functional relation used,
and the final rational coefficient vector.  Recomputing a very small
numerical residual does not address the missing functional identity.
The colored kernels developed here can organize additional relations,
but no identity annihilating those two residuals has been proved in
this article.

\subsection{Minimal coordinates for the first triple Stieltjes periods}

Theorem~\ref{thm:triple-jets} establishes a finite Tornheim-jet description
at every index.  The next question is whether, at rational shifts,
symmetries and character decompositions reduce the number of independent
formal jet coordinates.  Start with the equal-thirds constant in
\eqref{eq:triple-third} and the first indices $(1,0,0)$ and $(1,1,0)$.
The integer-value identity at $C=0$ determines
$\mathcal T(A,B,0)=\Li_A(X)\Li_B(Y)$, but it does not determine
$\partial_C\mathcal T(A,B,0)$; the latter contains a true logarithm of
the coupled summation variable $m+n$.

A concrete deliverable would be a finite presentation of these low jet
orders modulo the proved symmetries, conjugation, distribution, and
functional equations, together with an exact reduction of the triple
digamma expression in that presentation.  A nonzero formal coordinate
would indicate that more relations are needed.  It would not prove
arithmetic independence of the associated complex numbers.

\subsection{Three unequal integer frequencies}

With three affine frequencies, the constraint is
$p_1n_1+p_2n_2+p_3n_3=0$ and generally leaves a lattice of rank two.
In each sign region, divisibility conditions arise when one frequency
is eliminated.  Finite character filters can express those conditions;
the remaining denominators involve positive integer linear forms.
The natural question is whether a finite family of weighted colored
Tornheim kernels gives a canonical formula compatible with the local
contact polynomial of Theorem~\ref{dil:contact}.

The first target should keep all singular sets disjoint and prove the
ordinary spectral identity in an absolutely convergent half-space.
One must then show that the continued kernel gives the specified
fixed-coordinate finite part.  Merely solving the lattice equation is
insufficient when the local scales differ.  Four or more translated
factors would introduce higher-dimensional analogues and a larger
set of sign regions.

\subsection{Collision identities and nonlinear coordinate changes}

When two shifts coalesce, the disjoint-support product argument no
longer applies.  It is natural to compare three precisely defined
operations: the separated-shift formula as the separation tends to
zero, the finite part of the coincident product, and the regular
coefficient of a joint spectral continuation.  Exact difference
formulas between these operations could identify the complete local
counterterms without assuming they commute.

A related question replaces the affine coordinate $pt$ by
$pt+q_2t^2+q_3t^3+\cdots$.  For each fixed Stieltjes and derivative
order, only finitely many Taylor coefficients of this coordinate can
enter a distribution supported at the endpoint.  Derive that finite
law, recover $b_{n,k}(\log p)$ when the nonlinear coefficients vanish,
and check composition of coordinate changes exactly.  Such a law
would give a principled extension of the scale convention instead of
adding unrelated constants to special cases.

\subsection{Mixed harmonic Stieltjes data beyond the constant term}

The moving-pole theorem determines $R_{r,0}(a)$ and the ordinary
Stieltjes boundary $R_{0,k}(a)$.  The genuinely mixed coefficients
$R_{r,k}(a)$ with $r,k\ge1$ remain a distinct question.  A first target
is $R_{1,1}(a)$ at $a=1$, $1/2$, and a general positive rational shift.
Its convergent integral in \eqref{harm:eq:mixed-integrals} makes the
normalization exact; what is missing is an additional reduction.

One can also introduce an independent shift in the inner harmonic
denominators, or a color $\xi^n$ in the outer sum.  The beta/Gamma
summation class then changes.  Classify which choices still admit a
Gamma-quotient generating theorem and which require hypergeometric,
multiple-Lerch, or elliptic coordinates.  This is especially relevant
to avoiding the false inference that the uncolored half-shift formula
already evaluates the alternating S6 and S8 sums.

\subsection{Negative-balance identities and other summation theorems}

The polynomial identity \eqref{gs:residue-rigidity} is independent of
the comparison shift $c$.  A direct finite Bernoulli-polynomial proof,
separate from the analytic residue proof, would explain this rigidity
algebraically and could expose further parameter symmetries.
Different values of $M$ in \eqref{gs:FM} are related by explicit
Hurwitz terms; determine a convenient minimal normalization for tables
of the higher resonance jets.

The same program can be applied to balanced ${}_3F_2$ and other classical
summation formulas.  The issue is to identify the full asymptotic
coefficient of each possible pole before taking parameter derivatives.
Which summation theorems retain a finite Gamma-product answer after
these subtractions?  Which instead force multiple-zeta or higher
hypergeometric periods?  A useful result would classify one coherent
family with a complete mixed-jet theorem, rather than enumerate a few
unconnected low-order sums.

\subsection{The full polylogarithmic boundary function}

At rational $c$, equation~\eqref{gs:rational-Lerch} converts the comparison
term to finitely many colored polylogarithms and their order derivatives.
The other side, at $a=b=1/2$, already involves a complete elliptic
integral.  Determine a natural function class for the full
$z$-dependent mixed jets of \eqref{gs:polylog-generator}, and formulate
connection identities between $z=0$ and $z=1$ in that class.

The boundary constants proved here supply exact normalization data.
They do not imply that the entire family reduces to ordinary multiple
polylogarithms.  A first concrete study could treat the first two order
derivatives at the central-binomial point and compare their differential
equations with the explicitly known Stieltjes boundary constants.

\subsection{Exact symbolic generation and formal proof interfaces}

The rational-kernel and contact-polynomial generators are small enough
to serve as a first exact implementation target for ProveIt.  The
essential invariant for a rational kernel is the cancellation of the
sum of its simple-pole residues; losing it would replace a convergent
identity by separately divergent expressions.  The essential invariant
for a finite part is its named local coordinate and contact term.

A formal development could isolate reusable lemmas for locally normal
parameter families, Mellin subtraction, meromorphic coefficient
extraction, and products with disjoint singular supports.  These lemmas
would support every-order theorems without separate proofs for every
table entry.  Any computational search for additional constant
relations should preserve an exact ordered basket and return a
candidate status until an analytic derivation or an exact relation
certificate is supplied.
